\scanpage{93}
% L1–6
il est immédiat que la donnée de $J\tau=a$ ($a$ arbitraire dans $C$) détermine $\tau$ à une transformation arithmétique près, (car il détermine à une homothétie près le groupe des périodes engendré par $\omega_1$ et $\omega_2$, comme on vérifie très facilement).

% L6–11
Il suit aussitôt que $J$ est \deleted{\unsure{univoque}}\inserted{biunivoque} au voisinage de \deleted{tout} point du demi-plan supérieur qui n'est point double d'aucune transformation $\left(\begin{smallmatrix}a,b\\c,d\end{smallmatrix}\right)$ arithmétique, et en ces points seulement.

% L11–14
\deleted{\unsure{Dans ce cas|Donc de plus|Des autres}}\inserted{On} sait d'autre part qu'il y a exactement deux classes de points doubles, savoir $i$ $(i=-\frac1i)$ et $\frac12+i\frac{\sqrt3}{2}$.

% L14–17
Par suite, les seuls points singuliers \inserted{à distance finie} de la fonction inverse $\tau(\unsure{w|u})$ de $J(\tau)$ sont les deux points $J(i)$ et $J(\frac12+i\frac{\sqrt3}{2})$.

% L18–24
Le calcul donne respectivement $0$ et $1$, i.e $g_2=0$ et $g_3=0$ (car en partant de $g_2=0$ resp$^t$ $g_3=0$, on trouve bien $\frac{\omega_2}{\omega_1}=i$ resp$^t$ $\frac{\omega_2}{\omega_1}=\frac1{\unsure{2|3}}+i\frac{\sqrt3}{2}$, par un calcul d'intégrale \deleted{\unsure{a|e}}très simple). -- Ces propriétés sont à la base des théorèmes de Picard sur les fonctions entières.

% lower left
% archive sheet number: 92
